\includegraphics[width=1in]{media/image1.jpeg}

Arkusz zawiera informacje prawnie chronione do momentu rozpoczęcia egzaminu.

\begin{longtable}[]{@{}l@{}}
\textbf{WYPEŁNIA ZESPÓŁ NADZORUJĄCY}
\end{longtable}

\begin{enumerate}
\item Instrukcja dla zdającego.
\end{enumerate}

Zadanie 1. (0--1)

Liczba \(\sqrt{4}\) jest równa

A. \(1\)

B. \(2\)

Zadanie 2. (0--2)

Rozwiąż równanie \(x + 3 = 5\).

Zadanie 3.

Dana jest funkcja \(f(x) = 2x\) oraz jej wykres na rysunku.

\includegraphics[width=2in]{media/image7.png}

Zadanie 3.1. (0--1)

Podaj \(f(2)\).

Zadanie 3.2. (0--3)

Uzasadnij, że \(f\) jest rosnąca.

Zadanie 4. (0–1)

Tabela przedstawia dane.

\begin{longtable}[]{@{}ll@{}}
1 & 2 \\
3 & 4 \\
\end{longtable}

Podaj medianę.

Koniec

\textbf{MATEMATYKA}

\includegraphics{media/image9.png}
